\section{Technology comparison and the next affine--M\"obius gate}
\label{sec:tpc32-next-gate}

The exact identities above expose several familiar analytic forms, but none
of the cited technologies accepts the complete physical coefficient without
an additional argument.  We record the boundary to prevent a formal
rearrangement from being mistaken for a cancellation theorem.

\subsection{One M\"obius sign versus two affine signs}

Classical uniform bounds of Davenport type estimate a single oscillatory
M\"obius sum
\begin{equation}
 \sum_{d\asymp D}\mu(d)\omega(d)e(\theta d)
 \label{eq:tpc32-davenport-form}
\end{equation}
with arbitrary logarithmic saving when $D$ is a polynomial scale and
$\omega$ has controlled smoothness; see, for example,
\citet[Chapter~13]{IwaniecKowalski2004}.  Fixed-period characters and
Mellin phases can be included without changing the essential one-variable
form.  This is useful for a bare row coefficient, but opening a reflected
ultra leg changes the arithmetic problem.

Fix $\ell,j$ and a reflected divisor $s$ on primitive support.  The
condition
\begin{equation}
 s\mid \ell dj+h
 \label{eq:tpc32-reflected-divisibility}
\end{equation}
has a unique solution $d\equiv d_0\pmod s$, because
$(s,\ell j)=1$.  Write $d=d_0+st$.  The complementary quotient is
\begin{equation}
 U(t)=\frac{\ell j(d_0+st)+h}{s}
 =u_\star+\ell jt,
 \qquad
 u_\star=\frac{\ell jd_0+h}{s},
 \label{eq:tpc32-affine-quotient}
\end{equation}
and the two affine forms satisfy the nondegeneracy identity
\begin{equation}
 \boxed{su_\star-\ell jd_0=h\ne0.}
 \label{eq:tpc32-affine-determinant}
\end{equation}
After the exact sign fusion, a typical opened coefficient contains
\begin{equation}
 \mu(d_0+st)\mu(u_\star+\ell jt).
 \label{eq:tpc32-affine-mobius-correlation}
\end{equation}
This is a two-point affine M\"obius correlation, not the one-sign sum
\eqref{eq:tpc32-davenport-form}.

Logarithmically averaged two-point Chowla--Elliott theorems control fixed
nonproportional affine forms in an averaged regime \citep{Tao2016}.  The
physical packet here asks additionally for uniformity as $s,\ell,j$ and the
affine coefficients grow polynomially, for dyadic rather than logarithmic
averaging, for fixed residue factors, and for reassembly over all rows and
cutoffs.  We do not claim that the existing theorem supplies this package.

\subsection{Spectral and Kloosterman technology}

Bilinear and trilinear estimates for Kloosterman fractions can yield strong
cancellation once the variables, moduli, smooth weights, and coefficient
norms match the hypotheses of the relevant theorem
\citep{DukeFriedlanderIwaniec1997,BettinChandee2018}.  Modern exceptional
spectral large sieves provide further gains in carefully structured
families \citep{Pascadi2026LargeSieve}.  In the present shell, however, the
modulus is tied to reflected target divisors, the prime--M\"obius
coefficient depends on a moving factorization, and the generic pair mask is
still joint.  No parameter ledger in this paper verifies the hypotheses of
one of those stronger estimates.  They are possible technologies for a
future refined packet, not black boxes used in our theorems.

\subsection{A concrete next target}

The conditional transfer theorem reduces a high-$\beta$ closure of the raw
shell to the exact boundary $\chi\le1/400$.  A robust route with a fixed
margin is to
prove, for the actual physical amplitude and $C\asymp J$, a power-saving
estimate of the form
\begin{equation}
 \left|\sum_nA_C(n)\right|^2
 \ll X^{1/400-\eta}\sum_n|A_C(n)|^2
 \label{eq:tpc32-next-flatness-target}
\end{equation}
for some fixed $\eta>0$.  A viable proof would have to exploit information
absent from the abstract norm interface, such as:
\begin{enumerate}[label=\textup{(\alph*)}]
 \item cancellation in the affine correlation
 \eqref{eq:tpc32-affine-mobius-correlation} uniform over a controlled
 subfamily;
 \item a projective or operator-norm decomposition of the actual pair mask;
 \item an additional centering identity that annihilates the coherent
 determinant mode before absolute reassembly; or
 \item a decomposition into spectral/Kloosterman ranges with a complete
 conductor and norm ledger.
\end{enumerate}
Each option can be tested independently.  Failure of one does not refute
the others, and a sharp impossibility result for any proposed interface
would itself refine the route map.

For the full calibrated cutoff difference, the two drift legs also impose
the inherited saving $L^{-1}$; hence any final exponent must be intersected
with the fixed prime-source margin.  Even a proof of
\eqref{eq:tpc32-next-flatness-target} would close only this residual packet.
It would not by itself supply positivity, undo all earlier reductions, or
prove a prime-pair theorem.

\subsection{Scope of the advance}

The advance in this paper is that the determinant-frequency route no longer
ends at the vague instruction ``control the zero mode.''  It proves that
ordinary determinant frequencies already have the required size outside a
power-small exceptional set, identifies the remaining hard raw component of
the original calibrated residual with one special member of that set,
quantifies the exact additional flatness needed, and
shows that the current Hilbert-space data cannot force it.  The remaining
door is narrow, but it is still an open analytic estimate rather than a
claimed consequence of the existing large sieve.
